\section*{Bab \ref{ch:g11:vision-and-brain} --- Penglihatan dan Otak}

\begin{solution}{exo:g11:vision-and-brain:1}
Cahaya dari medan kirinya jatuh pada separuh kanan tiap retinanya;
serabut separuh kanan \omterm{def:g11:the-eye:parts}{mata} kirinya menyeberang di kiasmanya, serabut
\omterm{def:g11:the-eye:parts}{mata} kanannya tinggal di kanan; bersama-sama keduanya membentuk lajur
optik kanan, beralih di talamusnya dan mencapai \omterm{prop:g11:vision-and-brain:pathway}{korteks penglihatan
primer} kanannya.
\end{solution}

\begin{solution}{exo:g11:vision-and-brain:2}
Kedua saraf optiknya bertemu; serabut dari separuh hidung tiap retinanya
menyeberang ke sisi yang lain, serabut dari separuh pelipisnya tidak.
Di seberangnya tiap lajurnya membawa separuh medan penglihatannya.
\end{solution}

\begin{solution}{exo:g11:vision-and-brain:3}
Daerah warna (dunia terlihat kelabu), daerah gerakan (tanpa tanggapan
gerak), daerah wajah (wajah tidak dikenali).
\end{solution}

\begin{solution}{exo:g11:vision-and-brain:4}
Keplastisan: kemampuan hubungan neuron untuk dikuatkan, dilemahkan,
dibentuk atau dipangkas menurut pemakaiannya. Kurun genting: tahap awal
kehidupan yang selama itu pengalamannya membentuk korteksnya secara
menentukan dan tak terbalikkan.
\end{solution}

\begin{solution}{exo:g11:vision-and-brain:5}
Ia mengikat reseptor serotonin neuron dalam daerah penglihatannya,
sehingga mengganggu aliran isyaratnya dan menghasilkan halusinasi tanpa
cahaya yang sepadan.
\end{solution}

\begin{solution}{exo:g11:vision-and-brain:6}
Sesudah kiasmanya, di sebelah kiri: lajur optik kiri, peralihan
talamus kiri atau \omterm{prop:g11:vision-and-brain:pathway}{korteks penglihatan primer} kiri.
\end{solution}

\begin{solution}{exo:g11:vision-and-brain:7}
Serabut yang menyeberang berasal dari separuh hidung kedua retinanya,
yang melihat separuh luar (pelipis) medannya: pasiennya kehilangan
separuh kiri medan \omterm{def:g11:the-eye:parts}{mata} kirinya dan separuh kanan medan \omterm{def:g11:the-eye:parts}{mata} kanannya
--- yaitu penglihatan terowongan pada kedua sisinya.
\end{solution}

\begin{solution}{exo:g11:vision-and-brain:8}
Normal: $12 + 20 + 36 + 20 = 88\%$. Yang dikurangi: $2 + 3 + 5 + 30 = 40\%$,
yang hanya 10\% di antaranya secara utama.
\end{solution}

\begin{solution}{exo:g11:vision-and-brain:9}
Menutupnya memaksa korteksnya memakai masukan \omterm{def:g11:the-eye:parts}{mata} yang lemah,
sehingga hubungannya terjaga dan dikuatkan selama kurun gentingnya
alih-alih dipangkas. Pada orang dewasa kurunnya sudah lewat:
hubungannya tidak lagi menata ulang diri pada skala itu, sehingga \omterm{def:g11:the-eye:parts}{mata}
yang lemah tidak dapat merebut kembali wilayahnya.
\end{solution}

\begin{solution}{exo:g11:vision-and-brain:10}
\omterm{def:g11:the-eye:photoreceptors}{Sel kerucut} dan \omterm{def:g11:the-eye:photoreceptors}{opsinnya} normal; yang hilang adalah tahap korteks yang
menafsirkan nisbahnya. Rekaman listrik tanggapan retinanya terhadap
cahaya berwarna akan normal padanya dan tidak normal pada orang yang
buta warna; lagi pula cacatnya menyeluruh dan diperoleh, sedangkan
cacat mereka sebagian dan bawaan.
\end{solution}

\begin{solution}{exo:g11:vision-and-brain:11}
Bertahun-tahun latihan telah menguatkan dan memperbanyak hubungan yang
melayani jarinya, sehingga wilayah korteksnya membesar. Tanpa latihan,
daerahnya akan mengerut kembali sepanjang beberapa tahun, seperti
ditunjukkan penelitian jugglingnya.
\end{solution}

\begin{solution}{exo:g11:vision-and-brain:12}
Mula-mula ia melihat cahaya, warna dan bercak yang bergerak tetapi
tidak dapat mengenali bentuk, wajah atau kedalamannya: retinanya
bekerja, tetapi korteksnya tidak pernah membangun, selama kurun
gentingnya, hubungan yang menafsirkan isyaratnya. Ia akan belajar
memakai penglihatannya sebagian, dengan lambat, tetapi tidak pernah
selancar orang yang melihat sejak lahir.
\end{solution}

\begin{solution}{exo:g11:vision-and-brain:13}
Pada foveanya tiap kerucutnya punya \omterm{def:g10:cells-common-unit:cell}{sel} ganglionnya sendiri, sehingga
pusat medannya mengirimkan jauh lebih banyak serabut tiap derajat
daripada tepinya, tempat seratus \omterm{def:g11:the-eye:photoreceptors}{sel batang} berbagi satu; korteksnya
membagikan luasnya sebanding dengan serabutnya, jadi sebanding dengan
rinciannya, bukan dengan derajat medannya.
\end{solution}

\begin{solution}{exo:g11:vision-and-brain:14}
Kegiatan daerah itu adalah tanggapan warnanya, entah isyaratnya datang
dari \omterm{def:g11:the-eye:parts}{matanya} atau dari ingatannya; tanggapannya adalah apa yang
diperbuat daerahnya, bukan apa yang diantarkan \omterm{def:g11:the-eye:parts}{matanya}. Halusinasi
adalah kegiatan itu yang dimulai sebuah obat alih-alih oleh niat atau cahaya.
\end{solution}

\begin{solution}{exo:g11:vision-and-brain:15}
\omterm{def:g11:the-eye:parts}{Matanya} mengubah cahaya menjadi isyarat; melihat itu dikerjakan otaknya.
Seorang pasien dengan \omterm{def:g11:the-eye:parts}{mata} yang utuh dapat kehilangan warna, gerak atau
wajah akibat sebuah luka korteks; seekor anak kucing dengan \omterm{def:g11:the-eye:parts}{mata} yang
utuh menjadi buta secara fungsi pada salah satu \omterm{def:g11:the-eye:parts}{matanya} jika korteksnya
tidak dibiarkan menghubungkannya; dan sebuah obat yang tidak menyentuh
\omterm{def:g11:the-eye:parts}{mata} maupun struktur korteksnya membuat seseorang melihat apa yang
tidak ada. \omterm{def:g11:the-eye:parts}{Matanya} memasok; otaknya melihat.
\end{solution}

\begin{solution}{pb:g11:vision-and-brain:1}
\textbf{1.} \omterm{def:g11:the-eye:parts}{Mata} kirinya atau saraf optiknya, sebelum kiasmanya: hanya
serabut satu \omterm{def:g11:the-eye:parts}{matanya} yang terkena.

\textbf{2.} Lajur optik, talamus atau \omterm{prop:g11:vision-and-brain:pathway}{korteks penglihatan} kanannya:
cacat yang dibagi kedua \omterm{def:g11:the-eye:parts}{matanya} bagi separuh medan yang sama hanya
dapat terletak di tempat serabut kedua \omterm{def:g11:the-eye:parts}{matanya} bagi separuh itu telah
berkumpul bersama, yaitu sesudah kiasmanya.

\textbf{3.} Kiasmanya sendiri: serabut yang menyeberang, dari separuh
hidung kedua retinanya, yang melihat separuh pelipis medannya,
terpotong.

\textbf{4.} Daerah gerakannya, di seberang korteks primernya. Medan dan
ketajaman yang utuh menunjukkan bahwa korteks primernya dan segalanya
sebelum itu bekerja.

\textbf{5.} Retina B, C dan D utuh; retina A mungkin utuh, mungkin
tidak. Periksalah dengan menengok retinanya memakai oftalmoskop dan
merekam tanggapan listriknya terhadap kilatan.

\textbf{6.} Korteks kanannya menerima separuh kiri medannya.

\textbf{7.} $2^\circ$ pusatnya: \qty{1}{cm}; pita $2^\circ$ pada
$40^\circ$: \qty{0.1}{cm}. Nisbahnya 10.

\textbf{8.} Foveanya menjejalkan kerucutnya dengan rapat, yang
masing-masing punya \omterm{def:g10:cells-common-unit:cell}{sel} ganglionnya sendiri, sehingga derajat pusatnya
mengirimkan jauh lebih banyak serabut daripada tepinya; korteksnya
memberi derajat itu luas yang sebanding lebih besar.

\textbf{9.} Sebuah bintik buta kecil di pusat medan kirinya, yang
melumpuhkan bacaannya; jika lebih ke depan, sebuah daerah buta yang
lebih besar di tepi medan kirinya, yang kurang disadari.

\textbf{10.} Dengan menggerakkan \omterm{def:g11:the-eye:parts}{matanya}, pasiennya membawa tiap
katanya ke bagian retina yang korteksnya utuh, dengan memakai tepi
lukanya; lambat, tetapi mungkin.

\textbf{11.} $20 + 36 + 20 = 76\%$.

\textbf{12.} $2 + 3 + 5 + 30 = 40\%$, dan hanya 10\% yang \omterm{def:g11:the-eye:parts}{mata} kirinya
menjadi masukan utama atau setaranya.

\textbf{13.} Masukannya mencapai korteksnya tetapi hubungannya
dilemahkan dan dipangkas karena tidak terpakai; hubungan \omterm{def:g11:the-eye:parts}{mata} kanannya
memekar ke dalam ruangnya.

\textbf{14.} Bahwa penataan ulangnya terbatas pada kurun genting awal
kehidupannya; padanannya pada manusia adalah ambliopia, yang hanya
berkembang pada masa kanak-kanak dan tidak terobati sesudahnya.

\textbf{15.} Karena tidak ada \omterm{def:g11:the-eye:parts}{mata} yang diunggulkan, tidak ada
pengambilalihan: histogramnya tetap kira-kira setangkup tetapi dengan
neuron yang menanggapi lebih sedikit seluruhnya, dan penglihatan anak
kucingnya buruk pada kedua \omterm{def:g11:the-eye:parts}{matanya} --- pengurangan tanpa persaingan
melemahkan keduanya.

\textbf{16.} Otaknya menekan bayangan \omterm{def:g11:the-eye:parts}{mata} yang berbelok itu agar tidak
melihat ganda; karena tidak terpakai, hubungannya dipangkas selama
kurun gentingnya dan \omterm{def:g11:the-eye:parts}{matanya} menjadi lemah selamanya. Kacamata
membetulkan optiknya, tetapi \omterm{def:g11:the-eye:parts}{matanya} tetap tidak terpakai, sehingga
pemangkasannya berlanjut.

\textbf{17.} Menutup \omterm{def:g11:the-eye:parts}{mata} yang baik memaksa korteksnya memakai isyarat
\omterm{def:g11:the-eye:parts}{mata} yang lemah, sehingga hubungannya terjaga dan menguat.

\textbf{18.} Kurun gentingnya berakhir sekitar umur tujuh tahun;
sesudahnya korteksnya tidak lagi membagi ulang wilayahnya pada skala
itu, sehingga hubungan \omterm{def:g11:the-eye:parts}{mata} lemah yang hilang tidak dapat dibangun kembali.

\textbf{19.} Pada orang dewasa hubungannya sudah tegak dan mantap;
pengambilalihan oleh \omterm{def:g11:the-eye:parts}{mata} yang lain adalah gejala korteks yang sedang
berkembang, selama kurun gentingnya, bukan korteks yang sudah matang.

\textbf{20.} A: \omterm{def:g11:the-eye:parts}{matanya} atau saraf optiknya; C: kiasmanya; B: lajur
kanan, talamus atau korteks primernya; D: daerah gerakan di seberangnya.
Keplastisan korteksnya, yang terbatas pada kurun genting, membuat
penutupan \omterm{def:g11:the-eye:parts}{mata} anaknya berhasil pada umur empat tahun dan gagal pada
umur lima belas.
\end{solution}
